%%%%%%%%%%%%%%%%%%%%%%%%%%%%%%%%%%%%%%%%%%%%%%%%%%%%%%%%%%%%%%%%%%%%%%
% Lớp 12 | Một số yếu tố xác suất | Công thức xác suất toàn phần. Công thức Bayes
% Dạng 2D6?2-1: Công thức lý thuyết
% Số câu: 7  (sinh tự động bởi tach_cau.py; từ nay thầy sửa trực tiếp file này)
%%%%%%%%%%%%%%%%%%%%%%%%%%%%%%%%%%%%%%%%%%%%%%%%%%%%%%%%%%%%%%%%%%%%%%

%%%%%%%%%%%%%%%% TRẮC NGHIỆM 4 PHƯƠNG ÁN %%%%%%%%%%%%%%%%

% nguồn: Lop_12_Dai_so_Chuong_6
\begin{ex}%[2D6N2-1]%[2D6N2-1]%[ma:5d998794]
	Cho  $A$ và $B$ là hai biến cố bất kì, với $0<\mathrm{P}(B)<1$. Khi đó
	\choice 
		{$\mathrm{P}(B \mid A)=\dfrac{\mathrm{P}(A  B)}{\mathrm{P}(B)}$}
		 {\True $\mathrm{P}(A \mid B)=\dfrac{\mathrm{P}(A B)}{\mathrm{P}(B)}$}
		{$\mathrm{P}(A \mid B)=\dfrac{\mathrm{P}(B)}{\mathrm{P}(A  B)}$}
		{$\mathrm{P}(B  \mid A)=\dfrac{\mathrm{P}(A B)}{\mathrm{P}(\overline{B})}$}
	\loigiai{Cho $A$ và $B$ là hai biến cố bất  kì, với $0<\mathrm{P}(B)<1$. Khi đó $\mathrm{P}(A \mid B)=\dfrac{\mathrm{P}(A B)}{\mathrm{P}(B)}$.
	} 
\end{ex}

% nguồn: Lop_12_Dai_so_Chuong_6
\begin{ex}%[2D6N2-1]%[ma:75381532]
	Giả sử $A$ và $B$ là hai biến cố ngẫu nhiên thỏa mãn $\mathrm{P}\left(A\right)>0$ và $0<\mathrm{P}\left(B\right)<1$. Khẳng định nào  dưới đây \textbf{sai}?
	\choice 
		 {\True $\mathrm{P}\left(B|A\right)=\dfrac{\mathrm{P}\left(B\right)\mathrm{P}\left(A|B\right)}{\mathrm{P}\left(B\right)\mathrm{P}\left(A|B\right)+\mathrm{P}\left(A\right)\mathrm{P}\left(B|A\right)}$}
		 {$\mathrm{P}\left(B|A\right)=\dfrac{\mathrm{P}\left(B\right)\mathrm{P}\left(A|B\right)}{\mathrm{P}\left(B\right)\mathrm{P}\left(A|B\right)+\mathrm{P}\left(\overline{B}\right)\mathrm{P}\left(A|\overline{B}\right)}$}
		{$\mathrm{P}\left(B|A\right)=\dfrac{\mathrm{P}\left(B\right)\mathrm{P}\left(A|B\right)}{\mathrm{P}\left(A\right)}$} 
		{$\mathrm{P}\left(A\right)=\mathrm{P}\left(B\right)\mathrm{P}\left(A|B\right)+\mathrm{P}\left(\overline{B}\right)\mathrm{P}\left(A|\overline{B}\right)$} 
	\loigiai{ 
		Giả sử $A$ và $B$ là hai biến cố ngẫu nhiên thỏa mãn $\mathrm{P} \left(A\right)>0$ và $0<\mathrm{P} \left(B\right)<1$.  Khi đó, theo công thức Bayes, ta có
		\[\mathrm{P} \left(B|A\right)=\dfrac{\mathrm{P}  \left(B\right)\mathrm{P} \left(A|B\right)}{\mathrm{P}\left(B\right)\mathrm{P} \left(A|B\right)+\mathrm{P} \left(\overline{B}\right)\mathrm{P} \left(A|\overline{B}\right)}\]
		hay còn có thể  viết dưới dạng
	 \[\mathrm{P} \left(B|A\right)=\dfrac{\mathrm{P} \left(B\right)\mathrm{P} \left(A|B\right)}{\mathrm{P}\left(A\right)}.\]}
\end{ex}

% nguồn: Lop_12_Dai_so_Chuong_6
\begin{ex}%[2D6N2-1]%[ma:17a2f474]
	Cho $A$, $B$ là hai biến cố của một phép thử có không gian mẫu gồm hữu hạn kết  quả đồng khả năng, biết $0<\mathrm{P}(B)<1$. Phát biểu nào sau đây đúng?
	 \choice
		{$\mathrm{P}(A)=\mathrm{P}(B)\cdot  \mathrm{P}(A\mid B)+\mathrm{P}\left(\overline{B}\right)\cdot \mathrm{P}\left(\overline{B}\mid A\right)$}
		{$\mathrm{P}(A)=\mathrm{P}(B)\cdot \mathrm{P}(B\mid A)+\mathrm{P}\left(\overline{B}\right)\cdot  \mathrm{P}\left(\overline{B}\mid A\right)$}
		{$\mathrm{P}(A)=\mathrm{P}(B)\cdot  \mathrm{P}(B\mid A)+\mathrm{P}\left(\overline{B}\right)\cdot \mathrm{P}\left(A\mid\overline{B}\right)$}
		{\True $\mathrm{P}(A)=\mathrm{P}(B)\cdot \mathrm{P}(A\mid  B)+\mathrm{P}\left(\overline{B}\right)\cdot \mathrm{P}\left(A\mid\overline{B}\right)$}
	 \loigiai{
		Công thức đúng là  $\mathrm{P}(A)=\mathrm{P}(B)\cdot \mathrm{P}(A\mid B)+\mathrm{P}\left(\overline{B}\right)\cdot \mathrm{P}\left(A\mid\overline{B}\right)$.
	} 
\end{ex}

% nguồn: Lop_12_Dai_so_Chuong_6
\begin{ex}%[2D6N2-1]%[ma:d6b80a1d]
	Cho hai biến cố $M, N$ với $0 < \mathrm{P}(N) < 1$. Chọn khẳng định đúng trong các khẳng định sau  đây
	\choice 
		{\True $\mathrm{P}(M) = \mathrm{P}(\overline{N})\cdot\mathrm{P}(M|\overline{N}) + \mathrm{P}(N)\cdot\mathrm{P}(M|N)$} 
		 {$\mathrm{P}(M) = \mathrm{P}(N)\cdot\mathrm{P}(M|N) - \mathrm{P}(\overline{N})\cdot\mathrm{P}(M|\overline{N})$}
		{$\mathrm{P}(M)  = \mathrm{P}(\overline{N})\cdot\mathrm{P}(M|\overline{N}) - \mathrm{P}(N)\cdot\mathrm{P}(M|N)$}
		{$\mathrm{P}(M) = \mathrm{P}(N)\cdot\mathrm{P}(M|N) + \mathrm{P}(\overline{N})\cdot\mathrm{P}(M|\overline{N})$} 
	 \loigiai{
		Theo công thức xác  suất toàn phần, ta có
		$$ \mathrm{P}(M) = \mathrm{P}(M \cap N) + \mathrm{P}(M \cap \overline{N})  $$\\
		Theo công thức nhân xác  suất
		\begin{itemize}
			\item $\mathrm{P}(M  \cap N) = \mathrm{P}(N) \cdot \mathrm{P}(M|N)$ (vì $\mathrm{P}(N) > 0$ theo giả thiết).
			\item $\mathrm{P}(M \cap \overline{N})  = \mathrm{P}(\overline{N}) \cdot \mathrm{P}(M|\overline{N})$.
		\end{itemize}
		Vì  $0 < \mathrm{P}(N) < 1$, nên $\mathrm{P}(\overline{N}) = 1 - \mathrm{P}(N)$ cũng thỏa $0 < \mathrm{P}(\overline{N}) < 1$, do đó $\mathrm{P}(\overline{N}) > 0$.\\
		Thay  vào công thức xác suất toàn phần, ta được
		$$ \mathrm{P}(M) = \mathrm{P}(N)\mathrm{P}(M|N) + \mathrm{P}(\overline{N})\mathrm{P}(M|\overline{N}). $$ 
		 Công thức này có thể viết lại thứ tự các số hạng
		$$ \mathrm{P}(M) = \mathrm{P}(\overline{N})\mathrm{P}(M|\overline{N}) +  \mathrm{P}(N)\mathrm{P}(M|N). $$
	} 
\end{ex}

% nguồn: Lop_12_Dai_so_Chuong_6
\begin{ex}%[2D6N2-1]%[ma:0aff78d6]
	Cho hai biến cố $A$ và $B$ với  $0<\mathrm{P}\left(B\right)<1$. Khẳng định nào sau đây là đúng?
	 \choice
		{\True $\mathrm{P}\left(A\right)=\mathrm{P}\left(B\right)\cdot\mathrm{P}\left(A\mid  B\right)+\mathrm{P}\left(\overline{B}\right)\cdot\mathrm{P}\left(A\mid\overline{B}\right)$}
		{$\mathrm{P}\left(A\right)=\mathrm{P}\left(A\right)\cdot \mathrm{P}\left(A\mid B\right)+\mathrm{P}\left(\overline{A}\right)\cdot \mathrm{P}\left(A\mid  \overline{B}\right)$}
		 {$\mathrm{P}\left(A\right)=\mathrm{P}\left(B\right)\cdot\mathrm{P}\left(A\mid\overline{B}\right)+\mathrm{P}\left(\overline{B}\right)\cdot \mathrm{P}\left(A\mid B\right)$}
		{$\mathrm{P}\left(A\right)=\mathrm{P}\left(B\right)\cdot\mathrm{P}\left(A\mid B\right)-\mathrm{P}\left(\overline{B}\right)\cdot\mathrm{P}\left(A\mid \overline{B}\right)$} 
	\loigiai{ 
		Theo công thức xác suất toàn phần, ta có $\mathrm{P}\left(A\right)=\mathrm{P}\left(B\right)\cdot \mathrm{P}\left(A\mid B\right)+\mathrm{P}\left(\overline{B}\right)\cdot  \mathrm{P}\left(A\mid\overline{B}\right)$.
	 }
\end{ex}

% nguồn: Lop_12_Dai_so_Chuong_6
\begin{ex}%[2D6N2-1]%[ma:d6df8d7c]
	Cho  hai biến cố $A$ và $B$, với $0<\mathrm{P}(B)<1$. Đẳng thức nào sau đây đúng.
	\choice 
		 {$\mathrm{P}(A)=\mathrm{P}(A\mid B) \cdot \mathrm{P}(B)+\mathrm{P}\big(A\mid \overline{B}\big) \cdot \mathrm{P}(A)$}
		 {\True $\mathrm{P}(A)=\mathrm{P}(A\mid B) \cdot \mathrm{P}(B)+\mathrm{P}\big(A\mid \overline{B}\big) \cdot \mathrm{P}\big(\overline{B}\big)$}
		{$\mathrm{P}(A)=\mathrm{P}(A\mid B) \cdot \mathrm{P}(A)+\mathrm{P}\big(\overline{A} \mid B\big)  \cdot \mathrm{P}\big(\overline{A}\big)$}
		{$\mathrm{P}(B)=\mathrm{P}(A\mid B) \cdot \mathrm{P}(B)+\mathrm{P}\big(A\mid  \overline{B}\big) \cdot \mathrm{P}\big(\overline{B}\big)$}
	\loigiai{ 
		Cho hai biến  cố $A$ và $B$, với $0< P(B) < 1$ khi đó $\mathrm{P}(A)=\mathrm{P}(A\mid B) \cdot \mathrm{P}(B)+\mathrm{P}\big(A\mid \overline{B}\big) \cdot \mathrm{P}\big(\overline{B}\big)$.
	 }
\end{ex}

% nguồn: Lop_12_Dai_so_Chuong_6
\begin{ex}%[2D6N2-1]%[ma:7df9d7d7]
	Giả sử $A$ và $B$ là  hai biến cố ngẫu nhiên thỏa mãn $\mathrm{P}\left(A\right)>0$ và $0<\mathrm{P}\left(B\right)<1$. Khẳng định nào sau đây đúng?
	\choice 
		 {$\mathrm{P}\left(B|A\right)=\dfrac{\mathrm{P}\left(B\right)+\mathrm{P}\left(A|B\right)}{\mathrm{P}\left(B\right)\mathrm{P}\left(A|B\right)+\mathrm{P}\left(\overline{B}\right)\mathrm{P}\left(A|\overline{B}\right)}$}
		{$\mathrm{P}\left(B|A\right)=\dfrac{\mathrm{P}\left(B\right)\mathrm{P}\left(A|B\right)}{\mathrm{P}\left(B\right)\mathrm{P}\left(A|B\right)-\mathrm{P}\left(\overline{B}\right)\mathrm{P}\left(A|\overline{B}\right)}$} 
		 {$\mathrm{P}\left(B|A\right)=\dfrac{\mathrm{P}\left(B\right)\mathrm{P}\left(A|B\right)}{\mathrm{P}\left(B\right)\mathrm{P}\left(A|\overline{B}\right)+\mathrm{P}\left(\overline{B}\right)\mathrm{P}\left(A|B\right)}$}
		 {\True $\mathrm{P}\left(B|A\right)=\dfrac{\mathrm{P}\left(B\right)\mathrm{P}\left(A|B\right)}{\mathrm{P}\left(B\right)\mathrm{P}\left(A|B\right)+\mathrm{P}\left(\overline{B}\right)\mathrm{P}\left(A|\overline{B}\right)}$}
	\loigiai{ 
		Giả sử $A$ và $B$ là hai  biến cố ngẫu nhiên thỏa mãn $\mathrm{P}\left(A\right)>0$ và $0<\mathrm{P} \left(B\right)<1$ , khi đó ta có công thức Bayes
		\[\mathrm{P} \left(B|A\right)=\dfrac{\mathrm{P}  \left(B\right)\mathrm{P} \left(A|B\right)}{\mathrm{P}\left(B\right)\mathrm{P} \left(A|B\right)+\mathrm{P} \left(\overline{B}\right)\mathrm{P} \left(A|\overline{B}\right)}.\]
		 hay
	 \[\mathrm{P} \left(B|A\right)=\dfrac{\mathrm{P} \left(B\right)\mathrm{P} \left(A|B\right)}{\mathrm{P}\left(A\right)}.\]}
\end{ex}
